\documentclass{writtenassignmentchat}

%% ── Word count ───────────────────────────────────────────────────────────────
%% Word count needs to be updated manually. To count, use the command `texcount --sum main.tex`. Change the "main" to your main file name.)
\renewcommand{\wordcountvalue}{000}

%% ── Line numbers ─────────────────────────────────────────────────────────────
%% Line numbers are ON by default. To disable, uncomment the next line:
% \showlinenumbersfalse

%% ── Graphic path ─────────────────────────────────────────────────────────────
% Location of figures relative to the .tex file



\title{Written assignment template for Chat's courses}


\begin{document}

\maketitle
\begin{flushleft}
Version: 1.0 (\today)

This document demonstrates the \texttt{writtenassignmentchat} \LaTeX{} class. Here are a few things to keep in mind:
\end{flushleft}

\begin{enumerate}
  \item When the info in this file contradicts with the instructions in your assignment sheet, follow the assignment sheet.
  \item \textbf{Do not} add your name and student ID in the text. (The assignment is usually printed and graded with a pseudonym.)
  \item Use the writing software that you are fluent with to draft your work (or to collaborate with others, if applicable). In the end, if you struggle with \LaTeX{}, you can use the MS Word template in the file \verb|z_fallback_template.docx| instead.
  \item Some assignments may not need title, section, figures, or table. Don't use these features when they are not needed.
  \item If you have suggestions to improve this template, file an issue on the  \linkedfootnote{https://github.com/chatchavan/assignment_template/issues}{template Github repository}.
\end{enumerate}


\section{Files}

\begin{itemize}
  \item \verb|0-main.tex| - Your writing goes here.
  \item \verb|a-references.bib| - Reference entries in the bibtex format.
  \item \verb|apalike-doi.bst| - bibliographical entry style.
  \item \verb|figures| - The folder where you put images.
  \item \verb|README.tex| - The source file of the template instruction.
  \item \verb|README.pdf| - The template instruction.
  \item \verb|writtenassignmentchat.cls| - The \LaTeX{} class file.
  \item \verb|z_fallback_template.docx| - A similar template in MS Word that you can use as a fallback.
\end{itemize}

\section{Getting started}
To use this template, add your content to \verb|0-main.tex|.


\section{Citations}

\subsection{Citation style}

We use APA-like citations because the last name of the first author appears in the text. 
This saves readers' time. (No need to hunt down number in the reference section.)
Examples:
%
\begin{itemize}
  \item You can cite parenthetically at the end of a sentence with the \verb|\citep| command  ~\citep{smith2021}.
  \item The citation can be a subject of the sentence. It read as if you say to "author and colleague...". Use the \verb|\citet| command for this: \citet{brown2020} presented some findings.
  \item Be precise with your citations when possible by, e.g., pointing to specific page number or section number. Example \verb|\citep[p. 12]{smith2021}| generates \citep[p. 12]{smith2021}.
  \item Alternatively, you can also a prefix text before parenthetical citation with a command like\\ \verb|\citep[see p. 12 of][]{smith2021}|. This results in \citep[see p. 12 of][]{smith2021}. 
  \item To cite multiple items in one parenthesis, use \citep{smith2021, brown2020}. This template does not support citing multiple items with text prefix/suffix.
\end{itemize}
%

\subsection{Reference section}
Ensure that the reference information is as complete as possible. 
Whenever available, provide Digital Object Identifier (DOI).
The bibliography style is fixed to \texttt{apalike} in the class; it does not need to be
declared separately in the \texttt{.tex} file.



\section{Word Count}

Immediately before the reference list, the class prints the word count of the body text, excluding references. The value is set in the preamble with \verb|\renewcommand{\wordcountvalue}{N}|. 
For automatic counting, run \texttt{texcount -{}-sum 0-main.tex} and copy the result.

If you use Overleaf, see \linkedfootnote{https://docs.overleaf.com/writing-and-editing/using-word-count}{how to get the word count from the GUI}. If you cannot get the word count to work in the command line, don't spend too much time figuring it out. Just compile your work in PDF, copy-and-paste it to MS Word, and use its word count feature.

\section{URLs and Hyperlinks}

When you include a URL, make it as linked text with the full address
repeated in a footnote, so the link is accessible in both digital and printed
versions.
Use the custom command \verb|\linkedfootnote{url}{text}| to automatically create this.
Example: \linkedfootnote{https://www.latex-project.org}{the \LaTeX{} project website}.


\section{Elements other than the text}
Some assignments are text-only, in that case, this section is not relevant.

\subsection{Figures}

\autoref{fig:fullwidth} provides an example of a figure. 
When you use the command \verb|\includegraphics|, it first searches the subfolder \texttt{figures}.
If  you want to change the path, add the like \verb|\graphicspath{{your-path/}}| before \verb|\begin{document}|.
We recommend creating figure with exact physical size that it will appear on the page, such that you can just use the command \texttt{includegraphics} without any parameter in \LaTeX. In this way, you can ensure that the original figure has appropriate resolution.

\begin{figure}[ht]
  \centering
  
\includegraphics{example-fullwidth.pdf}
  \caption{Example figure that spans full width of the document.}
  \label{fig:fullwidth}
\end{figure}

\subsection{Table}

Table in \LaTeX{} can be tricky. For assignments, there's no need to spend too much time wrangle with \LaTeX{} syntax. Instead, you can create a table in other app (e.g., MS Word), export it to PDF, crop it, and include the table as image like \autoref{tab:tablefrompdf}. 


\begin{table}[ht]
  \centering
  \caption{Example table created in \LaTeX}
  \label{tab:tableinlatex}
  \begin{tabular}{llc}
    \toprule
    \textbf{Element} & \textbf{Setting} & \textbf{Value} \\
    \midrule
    Body text    & Font family   & Helvetica (helvet) \\
                 & Font size     & 11\,pt \\
                 & Line spacing  & 1.3$\times$ \\
    \midrule
    Section      & Font size     & 12\,pt \\
                 & Style         & Bold \\
    \midrule
    Subsection   & Font size     & 11\,pt \\
                 & Style         & Italic, underlined \\
    \bottomrule
  \end{tabular}
\end{table}

\begin{table}[ht]
  \centering
  \caption{Example table imported from a PDF file}
  \label{tab:tablefrompdf}
  \includegraphics{example-table-as-figure.pdf} 
\end{table}

\subsection{Line numbers}

Line numbers in the left margin can be disabled by adding the command \verb|\showlinenumbersfalse| before  \verb|\begin{document}|.

\section{Stylistic specifications}
If you use the template, the following settings are already applied. They are documented here for reference purposes.

\subsection{Font style}

Body text is set in \textbf{Helvetica} at \textbf{11\,pt} with a
line spacing of \textbf{1.3} lines.
Section headings are \textbf{12\,pt bold}.
Subsection headings are \textbf{11\,pt italic and underlined}. 

Do not indent the first paragraph of each section. Do indent subsequent paragraphs. Do not use ligatures.

\subsection{Margins}

The page margins are \textbf{1\,cm top}, \textbf{1.5 cm bottom} and \textbf{2\,cm left and
right}. No header is printed; the footer
area is included in the bottom margin calculation.

\section{Acknowledgement}
This template was scaffoled with multiple iterations with Claude Sonnet 4.6, plus a few hours of manual tweaks.

\insertbibliography{a-references}

\end{document}
